\documentclass[10pt,twocolumn]{article}

% --- PACKAGES ---
\usepackage[utf8]{inputenc} % Input encoding
\usepackage[T1]{fontenc}    % Font encoding
\usepackage[english]{babel} % Language
\usepackage{graphicx}       % For including images
\usepackage{amsmath}        % For math formulas
\usepackage{hyperref}       % For clickable links (references, URLs)
\hypersetup{colorlinks=true, linkcolor=blue, citecolor=blue, urlcolor=blue}
\usepackage{geometry}       % For page margins
\geometry{margin=1.5cm}     % Set margins (adjust as needed)
\usepackage{titlesec}       % For customizing section titles
\usepackage{caption}        % For customizing captions
\usepackage{float}          % For better figure placement (use [H] option)
\usepackage{booktabs}       % For professional-looking tables
\usepackage{tabularx}       % For tables with adjusted column widths
\usepackage{enumitem}       % For customizing lists
\usepackage{times}          % Use Times font for a classic report look (optional)
\usepackage{url}            % Handles URLs better, especially with special chars

% --- FORMATTING ---
% Section/Subsection Titles
\titleformat{\section}{\large\bfseries}{\thesection}{1em}{}
\titleformat{\subsection}{\normalsize\bfseries}{\thesubsection}{1em}{}
\titleformat{\subsubsection}{\small\bfseries}{\thesubsubsection}{1em}{} % Added for more structure if needed

% Figure/Table Captions
\captionsetup[figure]{labelfont=bf, name=Figure, labelsep=period}
\captionsetup[table]{labelfont=bf, name=Table, labelsep=period}

% --- DOCUMENT ---
\title{\textbf{Sentiment Analysis of Tweets: A Comparative Study of Lexicon-Based, SVM, and LSTM Approaches using Cloud Data Management}} % Updated Title
\author{Amine Sekkat,  \\ Master in Computer Science - Data Science \\ % Add University/Affiliation if desired
        \texttt{your.email@example.com}} % Optional: Add email % Use \texttt for email
\date{\today} % Use current date

\begin{document}

\maketitle
\thispagestyle{empty} % No page number on title page

\begin{abstract}
Sentiment analysis of social media text, particularly tweets, presents unique challenges due to inherent noise, brevity, and informal language. This report presents a comprehensive study comparing three distinct approaches for classifying tweet sentiment (positive, negative, neutral): a lexicon-based method leveraging TextBlob's predefined sentiment scores, a classical machine learning model using Support Vector Machines (SVM) with TF-IDF features, and a deep learning model employing Long Short-Term Memory (LSTM) networks. The study utilizes a publicly available labeled Twitter dataset, managed within a cloud-based MongoDB Atlas database to facilitate scalability and flexible access. We evaluate these methods rigorously using standard metrics including accuracy, precision, recall, and F1-score, analyzing both overall performance and per-class results. Furthermore, we investigate model behavior through specific analyses like polarity distributions, feature importance, and error patterns based on text length. The objective is to provide a clear understanding of the strengths, weaknesses, and practical considerations of each approach in the context of real-world tweet sentiment analysis, while also justifying the chosen data management strategy.
\end{abstract}

\section{Introduction}
\label{sec:introduction}

In the contemporary digital era, social media platforms like Twitter have become vast repositories of public opinion and real-time reactions to global events. Extracting meaningful insights from this deluge of user-generated text is crucial for businesses, policymakers, and researchers alike. Sentiment analysis, a subfield of Natural Language Processing (NLP), aims to automatically identify and categorize opinions expressed in text, determining whether the writer's attitude towards a particular topic or entity is positive, negative, or neutral. Applications range widely, encompassing brand monitoring, customer feedback analysis, political campaign tracking, public health surveillance, and market research \cite{ref:feldman2013}.

However, performing sentiment analysis on tweets is notoriously difficult. Tweets are characterized by:
\begin{itemize}[leftmargin=*, itemsep=0pt, topsep=3pt]
    \item \textbf{Brevity:} Strict character limits (historically 140, now 280) force concise, often incomplete sentences.
    \item \textbf{Informality:} Rampant use of slang, abbreviations, misspellings, and non-standard grammar.
    \item \textbf{Noise:} Presence of URLs, hashtags, user mentions, and emojis that can obscure or modify sentiment.
    \item \textbf{Context Dependency:} Sentiment can be highly dependent on the context of the conversation or event. % Fixed leading character
    \item \textbf{Implicit Sentiment:} Sarcasm, irony, and humor often convey sentiment opposite to the literal meaning of the words used. % Fixed leading character
\end{itemize}
These factors pose significant challenges for traditional NLP techniques.

This project undertakes an experimental comparison of three widely used sentiment analysis methodologies applied to a standard Twitter dataset:
\begin{enumerate}[label=\arabic*), leftmargin=*, itemsep=0pt, topsep=3pt]
    \item \textbf{Lexicon-Based Approach:} Utilizes TextBlob, which relies on predefined dictionaries mapping words to polarity scores. It's simple but potentially limited by vocabulary coverage and context insensitivity.
    \item \textbf{Support Vector Machine (SVM):} A robust classical machine learning algorithm known for effectiveness in high-dimensional text classification tasks. We employ TF-IDF (Term Frequency-Inverse Document Frequency) features to represent the tweets numerically.
    \item \textbf{Long Short-Term Memory (LSTM):} A type of Recurrent Neural Network (RNN) specifically designed to handle sequential data and capture long-range dependencies, potentially offering advantages in understanding tweet context.
\end{enumerate}

By implementing and evaluating these diverse techniques, we aim to understand their relative performance characteristics, computational requirements, and suitability for the specific task of tweet sentiment classification. Furthermore, we address the practical aspect of data management by utilizing a cloud-based NoSQL database (MongoDB Atlas) and justify this choice in the context of the project's needs. The insights gained can inform the selection of appropriate methods for similar real-world applications.

\section{Related Work}
\label{sec:related_work}
Sentiment analysis on Twitter has been an active research area for over a decade. Numerous approaches, ranging from simple rule-based systems to complex deep learning architectures, have been proposed. This section briefly reviews relevant works categorized by methodology, similar to the approaches explored in this project, and positions our work within this landscape. (Aiming for ~1.5 pages).

\subsection{Lexicon-Based Methods}
Lexicon-based approaches are among the earliest and simplest methods. They rely on sentiment lexicons (dictionaries) where words are pre-assigned polarity scores (e.g., positive, negative, neutral) and sometimes subjectivity scores. The sentiment of a text is often determined by aggregating the scores of its constituent words.

Examples include SentiWordNet \cite{ref:sentiwordnet}, VADER (Valence Aware Dictionary and sEntiment Reasoner) \cite{ref:vader}, and the lexicons used by libraries like TextBlob \cite{ref:textblob}. Catelli et al. (2023) \cite{ref:catelli2023} demonstrated a recent application, using a custom Italian lexicon for analyzing COVID-19 related tweets. Their work highlights the domain-specificity often required for high performance, as general-purpose lexicons may miss domain-specific slang or meanings.

**Strengths:** These methods are generally fast, computationally inexpensive, require no training data (unsupervised), and are easily interpretable.
**Weaknesses:** Performance heavily depends on lexicon quality and coverage. They struggle with context, negation handling (though some like VADER attempt this), sarcasm, irony, and domain-specific language not present in the lexicon. They often perform poorly compared to supervised machine learning methods.

\subsection{Classical Machine Learning Approaches}
Supervised machine learning algorithms have been widely applied to sentiment analysis, treating it as a text classification problem. These methods require labeled training data. Common algorithms include Naive Bayes (NB), Logistic Regression (LR), and Support Vector Machines (SVM).

A crucial step is feature engineering: converting raw text into numerical feature vectors. Popular techniques include Bag-of-Words (BoW), Term Frequency (TF), Term Frequency-Inverse Document Frequency (TF-IDF), and N-grams. Ahmad et al. (2017) \cite{ref:ahmad2017} provide a typical example using SVM with TF-IDF features for tweet classification, achieving respectable results. Pang and Lee's seminal work \cite{ref:pang_lee_sentiment} also established strong baselines using these techniques on movie reviews.

**Strengths:** Generally outperform lexicon-based methods when sufficient labeled data is available. SVMs, in particular, handle high-dimensional sparse data (like TF-IDF vectors) effectively. Relatively well-understood algorithms with available libraries (e.g., Scikit-learn).
**Weaknesses:** Require labeled training data. Feature engineering choices significantly impact performance. Models like SVM can be sensitive to parameter tuning (e.g., kernel, C parameter). May still struggle with complex linguistic phenomena like sarcasm and long-range context compared to deep learning.

\subsection{Deep Learning Approaches}
With the rise of deep learning, various neural network architectures have been applied to sentiment analysis, often achieving state-of-the-art results. Recurrent Neural Networks (RNNs), particularly LSTMs \cite{ref:lstm_hochreiter} and Gated Recurrent Units (GRUs) \cite{ref:gru_cho}, are well-suited for sequential data like text, as they can capture contextual information over time. Convolutional Neural Networks (CNNs) \cite{ref:cnn_kim} have also been used, often capturing local N-gram patterns effectively.

Roy et al. (2024) \cite{ref:roy2024} specifically propose an LSTM-based method for hateful sentiment detection in tweets, comparing it favorably against classical ML models and even BERT. Their work highlights the potential of LSTMs for real-time analysis. Attention mechanisms \cite{ref:attention_bahdanau} are often combined with RNNs or used in Transformer models \cite{ref:transformer_vaswani} like BERT (Bidirectional Encoder Representations from Transformers) \cite{ref:bert_devlin} and its variants (RoBERTa, ALBERT, etc.). These Transformer models, pre-trained on massive text corpora, currently represent the state-of-the-art for many NLP tasks, including sentiment analysis, as they capture deep contextual relationships.

**Strengths:** Can automatically learn relevant features from raw text (often using word embeddings like Word2Vec \cite{ref:word2vec} or GloVe \cite{ref:glove}, or contextual embeddings from Transformers). Capture complex patterns, context, and sequential dependencies effectively. Often achieve the highest accuracy.
**Weaknesses:** Require large amounts of labeled training data. Computationally expensive to train, often requiring GPUs. Can be harder to interpret (less explainable) than classical models. Sensitive to hyperparameter tuning.

\subsection{Comparison and Positioning of Our Work}
This project sits at the intersection of these approaches by implementing and directly comparing a representative method from each category: TextBlob (Lexicon), Linear SVM (Classical ML), and LSTM (Deep Learning). Unlike studies focusing on a single method, our goal is a comparative evaluation on the same dataset and task (3-class sentiment classification).

Our work is similar in spirit to comparative studies like that of Roy et al. \cite{ref:roy2024}, but focuses on general positive/negative/neutral sentiment rather than just hateful content, and includes a direct lexicon-based comparison. We use standard TF-IDF features for SVM and word embeddings followed by LSTM layers for the deep learning approach. A key differentiator is our use and discussion of a cloud-based NoSQL database (MongoDB Atlas) for data management, an aspect often overlooked in purely model-focused sentiment analysis papers. While we do not propose a novel architecture, the practical comparison and the integration of data management considerations provide valuable insights for practitioners choosing a method for tweet analysis. Our work aims to reproduce standard implementations and evaluate their trade-offs in a clear experimental setting.

% --- Aim for ~1.5 pages for Related Work ---

\section{Problem Formulation}
\label{sec:problem_formulation}

Let $D = \{(x_i, y_i)\}_{i=1}^{n}$ be the dataset, where $n$ is the total number of tweets. Each instance $i$ consists of:
\begin{itemize}[itemsep=0pt, topsep=3pt]
    \item $x_i$: The raw text content of the $i$-th tweet.
    \item $y_i$: The true sentiment label associated with tweet $x_i$.
\end{itemize}
The set of possible sentiment labels is $Y = \{\text{positive, neutral, negative}\}$. % Used \text{} for robustness

The objective of this project is to develop and evaluate three distinct classification models, denoted $f_{lexicon}$, $f_{SVM}$, and $f_{LSTM}$. Each model $f$ aims to learn a mapping from the input tweet text $x$ to a predicted sentiment label $\hat{y} \in Y$:
$$ \hat{y} = f(x) $$
The goal is to find models $f$ that generalize well to unseen tweets, maximizing agreement between the predicted label $\hat{y}$ and the true label $y$ on a held-out test set $D_{test} \subset D$. Performance is measured using standard classification metrics such as Accuracy, Precision, Recall, and F1-score, calculated typically on $D_{test}$. We compare the performance of $f_{lexicon}$, $f_{SVM}$, and $f_{LSTM}$ based on these metrics.

\section{Proposed Solution}
\label{sec:solution}

We implement and compare three distinct approaches for tweet sentiment analysis. All models utilize data stored and retrieved from a MongoDB Atlas cloud database instance.

\subsection{Data Management: MongoDB Atlas}
\label{subsec:data_management}

Choosing an appropriate data management strategy is crucial, especially when dealing with potentially large volumes of semi-structured data like tweets. For this project, we opted to store the dataset in **MongoDB Atlas**, a cloud-based NoSQL document database service, instead of relying solely on local flat files (like CSV or Parquet).

**Justification:**
\begin{itemize}[leftmargin=*, itemsep=0pt, topsep=3pt]
    \item \textbf{Scalability:} Cloud databases like MongoDB Atlas are designed for horizontal scalability. While our current dataset might fit in memory, using a cloud DB demonstrates a solution prepared for larger, real-world data volumes where single-machine processing becomes infeasible.
    \item \textbf{Flexibility:} MongoDB's document model (BSON) handles semi-structured data like tweets well. It doesn't require a rigid predefined schema like relational databases, making it easier to accommodate variations in tweet structure or add new fields later (e.g., storing model predictions, confidence scores, or extracted entities alongside the original tweet).
    \item \textbf{Accessibility \& Collaboration:} A cloud database allows easy access for potentially distributed team members or deployment in cloud environments, decoupling the data storage from the local machine running the analysis code.
    \item \textbf{Managed Service:} Atlas handles administrative tasks like backups, patching, and monitoring, allowing focus on the analysis rather than database maintenance. % Fixed item marker
    \item \textbf{Efficient Querying:} MongoDB provides rich querying capabilities, including indexing (e.g., on the `sentiment` field to quickly retrieve tweets of a specific class) and projection (used in our scripts to fetch only necessary fields like `text` and `sentiment`, reducing network transfer and memory usage). % Fixed item marker
    \item \textbf{Aggregation Framework:} As demonstrated in our supplementary analysis (Section \ref{subsec:exp_aggregation}), MongoDB's aggregation pipeline allows performing data analysis (like calculating counts and averages per sentiment) directly within the database, which can be more efficient than transferring raw data for simple analytics. % Fixed item marker
\end{itemize}

While formats like Parquet offer columnar efficiency benefits, MongoDB provided a balanced solution offering scalability, flexibility, and powerful querying suitable for the iterative nature of this experimental project. Data is fetched from MongoDB into Pandas DataFrames within each Python script using the \texttt{pymongo} library at the start of the workflow.

\subsection{Approach 1: Lexicon-Based Model (TextBlob)}
\label{subsec:sol_lexicon}

This approach utilizes the pre-built sentiment analysis functionality of the TextBlob Python library \cite{ref:textblob}. TextBlob provides a simple API based on NLTK and Pattern libraries, including a sentiment property that returns polarity and subjectivity scores.

\textbf{Preprocessing:} Before analysis, minimal preprocessing is applied to the tweet text (or \texttt{selected\_text} in some experiments): converting to lowercase, removing URLs, user mentions, and hashtags. Punctuation removal was also tested but sometimes hinders TextBlob's performance.
% \texttt{preprocess\_text(text)} function used: [Include brief description or pseudocode if desired]

\textbf{Sentiment Calculation:} For a given tweet text $x$, TextBlob calculates:
\begin{itemize}[itemsep=0pt, topsep=3pt]
    \item Polarity ($p(x)$): A float within the range [-1.0, 1.0], where -1.0 is very negative and 1.0 is very positive.
    \item Subjectivity ($s(x)$): A float within the range [0.0, 1.0], where 0.0 is very objective and 1.0 is very subjective.
\end{itemize}
These scores are based on an internal lexicon derived primarily from the Pattern library's resources.

\textbf{Classification Rule:} We map the polarity score $p(x)$ to one of the three sentiment labels using predefined thresholds ($\theta_{pos} = 0.1$, $\theta_{neg} = -0.1$, based on experimentation; other thresholds like 0.25 were also tested):
$$ \hat{y}_{lexicon} = f_{lexicon}(x) =
\begin{cases}
    \text{positive} & \text{if } p(x) > \theta_{pos} \\
    \text{negative} & \text{if } p(x) < \theta_{neg} \\
    \text{neutral}  & \text{otherwise}
\end{cases}
$$
% Pseudocode can be added if needed

\textbf{Implementation:} This was implemented using the \texttt{TextBlob} class from the library. The \texttt{get\_sentiment} function encapsulates the threshold logic. Evaluation involves applying this function to the test set texts fetched from MongoDB.

\subsection{Approach 2: SVM Classifier (TF-IDF)}
\label{subsec:sol_svm}

This approach uses a classical supervised machine learning pipeline with a Support Vector Machine classifier.

\textbf{Preprocessing:} Raw tweet text fetched from MongoDB undergoes cleaning: removing URLs, mentions, hashtags, punctuation (except perhaps important ones if needed), converting to lowercase, and stripping extra whitespace.
% \texttt{clean\_text\_svm(text)} function used: [Include brief description or pseudocode if desired]

\textbf{Feature Extraction (TF-IDF):} The cleaned text data needs to be converted into numerical vectors suitable for the SVM. We use the Term Frequency-Inverse Document Frequency (TF-IDF) representation:
\begin{itemize}[itemsep=0pt, topsep=3pt]
    \item A vocabulary of the top $V=5000$ most frequent unigrams and bigrams (N-grams of range (1, 2)) is built from the training corpus, ignoring common English stop words.
    \item Each tweet $x_i$ is represented as a $V$-dimensional vector $\mathbf{v}_i$, where each component $v_{ij}$ corresponds to the TF-IDF weight of the $j$-th N-gram in the vocabulary for tweet $x_i$.
    \item TF-IDF assigns higher weights to terms that are frequent in a document but rare across the entire corpus, highlighting potentially discriminative words.
\end{itemize}
This is implemented using \texttt{TfidfVectorizer} from Scikit-learn.

\textbf{Classification Model (Linear SVM):} We employ a Linear Support Vector Machine (LinearSVC) \cite{ref:svm_cortes}. SVMs aim to find an optimal hyperplane that separates data points belonging to different classes in the high-dimensional feature space created by TF-IDF.
\begin{itemize}[itemsep=0pt, topsep=3pt]
    \item For multi-class classification (positive, negative, neutral), Scikit-learn's \texttt{LinearSVC} typically uses a "one-vs-rest" (OvR) strategy by default. It trains a separate binary classifier for each class against all other classes.
    \item Key parameter: $C=0.5$ (regularization parameter, controls trade-off between margin maximization and misclassification penalty) was used based on preliminary tests (or specify if default was used).
\end{itemize}
The model is trained on the TF-IDF vectors of the training set (\texttt{X\_train\_vec}, \texttt{y\_train}). % Escaped underscores

\textbf{Implementation:} Implemented using Scikit-learn's \texttt{TfidfVectorizer} and \texttt{LinearSVC}. The script \texttt{svm\_classifier.py} loads data, performs the split, fits the vectorizer and model, makes predictions on the test set, evaluates, and saves the model/vectorizer using \texttt{joblib}. % Escaped underscores

\subsection{Approach 3: LSTM Classifier}
\label{subsec:sol_lstm}

This approach utilizes a Long Short-Term Memory (LSTM) network, a type of RNN suitable for sequence modeling. The implementation described here focuses on binary classification (positive vs. negative) as per the provided \texttt{data\_preprocessing.py} script, although the architecture can be adapted for 3 classes. % Escaped underscore

\textbf{Preprocessing:} Similar initial text cleaning (lowercase, remove URLs, mentions, punctuation) is applied as in the SVM approach. Crucially, tweets labeled as 'neutral' are filtered out, leaving only 'positive' and 'negative' samples for this specific model implementation.
% \texttt{clean\_text\_lstm(text)} function used: [Include brief description or pseudocode if desired]

\textbf{Tokenization and Padding:}
\begin{itemize}[itemsep=0pt, topsep=3pt]
    \item A vocabulary of size $V=5000$ is built from the training texts using \texttt{Tokenizer} from Keras. Words not in the vocabulary are mapped to a special '<OOV>' token.
    \item Each tweet is converted into a sequence of integer indices corresponding to the words in the vocabulary.
    \item All sequences are padded or truncated to a fixed maximum length ($L=28$) using post-padding. This results in input tensors of shape \texttt{(n\_samples, max\_len)}. % Escaped underscores
\end{itemize}

\textbf{Label Encoding:} The 'positive'/'negative' labels are converted into a one-hot encoded format (e.g., [0, 1] for positive, [1, 0] for negative) using \texttt{pd.get\_dummies}, resulting in target tensors of shape \texttt{(n\_samples, 2)}. % Escaped underscore

\textbf{LSTM Model Architecture:} The Keras Sequential model consists of:
\begin{enumerate}[label=\alph*), itemsep=0pt, topsep=3pt]
    \item \textbf{Embedding Layer:} Maps each integer index in the input sequence to a dense vector representation of dimension $E=128$. It learns word embeddings from scratch during training. \texttt{input\_dim=V=5000}, \texttt{output\_dim=E=128}, \texttt{input\_length=L=28}. % Escaped underscores
    \item \textbf{SpatialDropout1D Layer:} A form of dropout ($rate=0.4$) that drops entire 1D feature maps, helping to prevent overfitting in convolutional/recurrent layers dealing with sequence data.
    \item \textbf{LSTM Layer:} The core recurrent layer with $N_{lstm}=196$ units. It processes the sequence of embeddings, maintaining an internal state to capture temporal dependencies. Dropout ($rate=0.2$) and recurrent dropout ($rate=0.2$) are applied for regularization. \texttt{return\_sequences=False} as it's the last recurrent layer feeding into a Dense layer. % Escaped underscore
    \item \textbf{Dense Output Layer:} A fully connected layer with 2 neurons (one for each class) and a `softmax` activation function, outputting probabilities for the 'negative' and 'positive' classes.
\end{enumerate}

\textbf{Training:} The model is compiled using the `categorical_crossentropy` loss function (appropriate for one-hot encoded labels with softmax output) and the `Adam` optimizer. It is trained for a specified number of epochs (e.g., 10 epochs in the script) using a batch size of 32. A portion of the training data (10\%) is used as a validation set during training.

\textbf{Implementation:} Implemented using TensorFlow/Keras. The \texttt{data\_preprocessing.py} script handles data loading from MongoDB, cleaning, filtering, tokenizing, padding, and splitting. The \texttt{train\_lstm.py} script defines the model architecture, compiles it, trains it, evaluates it on the test set, and saves the model (\texttt{.keras}) and tokenizer (\texttt{.pickle}). % Escaped underscores

\section{Experiments}
\label{sec:experiments}

This section details the experimental setup, the dataset used, evaluation metrics, the results obtained for each model, and a discussion comparing their performance.

\subsection{Experimental Environment}
\label{subsec:exp_env}

All experiments were conducted on the following setup:
\begin{itemize}[itemsep=0pt, topsep=3pt]
    \item \textbf{Operating System:} Windows 10 [Or specify your OS]
    \item \textbf{Processor (CPU):} [Specify your CPU, e.g., Intel Core i7-8700K @ 3.70GHz]
    \item \textbf{Memory (RAM):} [Specify RAM, e.g., 16 GB]
    \item \textbf{GPU:} No GPU acceleration was utilized for these experiments. [Or specify GPU if used, e.g., NVIDIA GeForce GTX 1080 Ti]
    \item \textbf{Key Software:} Python 3.10 [Or your version], Pandas [Version], Scikit-learn [Version], TensorFlow [Version] (for LSTM), TextBlob [Version], Pymongo [Version], Joblib [Version], Matplotlib [Version], Seaborn [Version].
    \item \textbf{Database:} MongoDB Atlas cluster (e.g., M0 free tier or specify).
\end{itemize}
The code was executed within [Specify environment, e.g., VS Code with Python extension, Jupyter Notebooks].

\subsection{Dataset Description}
\label{subsec:exp_dataset}

The dataset used in this project is derived from a publicly available collection of tweets labeled for sentiment. It was initially provided as a CSV file (\texttt{tweet.csv}) containing columns such as \texttt{textID}, \texttt{text} (the full tweet), \texttt{selected\_text} (the text span supporting the sentiment), and \texttt{sentiment}. % Escaped underscores

For this project, the data was ingested into a MongoDB Atlas database (\texttt{sentiment\_project} database, \texttt{tweets} collection). Data fetching operations retrieved only the necessary fields (\texttt{text}, \texttt{sentiment}) for model training and evaluation, leveraging MongoDB's projection capabilities. % Escaped underscore

**Preprocessing and Filtering:** Before splitting, records with missing `text` or `sentiment` values were dropped. Only records with `sentiment` labels explicitly marked as 'positive', 'negative', or 'neutral' were retained for analysis. Note that the LSTM implementation, as described in Section \ref{subsec:sol_lstm}, further filtered out 'neutral' tweets, training only on positive and negative examples. The SVM and Lexicon models were evaluated on all three classes.

**Data Split:** The cleaned dataset was consistently split into training and testing sets using an 80\%/20\% ratio using \texttt{train\_test\_split} from Scikit-learn (\texttt{random\_state=42} for reproducibility). Stratification based on the sentiment label (\texttt{stratify=y}) was used to ensure similar class distributions in both the training and testing sets. % Escaped underscores

**Dataset Statistics (Post-Cleaning, Pre-Split):**
\begin{itemize}[itemsep=0pt, topsep=3pt]
    \item Total number of tweets: [Insert Number, e.g., 27,480]
    \item Distribution:
        \begin{itemize}
            \item Positive: [Insert Count/Percentage, e.g., 8,582 (31.2\%)]
            \item Negative: [Insert Count/Percentage, e.g., 7,781 (28.3\%)]
            \item Neutral: [Insert Count/Percentage, e.g., 11,117 (40.5\%)]
        \end{itemize}
    \item Average text length: [Insert Average Length, e.g., approx. 69 characters based on aggregation]
\end{itemize}
*(Note: Obtain these numbers by running the exploration notebook or the aggregation script on your final cleaned dataset loaded from MongoDB).*

\begin{figure}[H] % Use [H] from float package for "here if possible" placement
    \centering
    % \includegraphics[width=0.9\columnwidth]{plots/sentiment_distribution.png} % Replace placeholder
    \framebox[0.9\columnwidth][c]{[Figure Placeholder: Dataset Sentiment Distribution]} % Placeholder Box
    \caption{Distribution of sentiment classes in the cleaned dataset.}
    \label{fig:sentiment_dist}
\end{figure}

\begin{figure}[H]
    \centering
    % \includegraphics[width=0.9\columnwidth]{plots/avg_length_by_sentiment.png} % Replace placeholder
    \framebox[0.9\columnwidth][c]{[Figure Placeholder: Average Text Length by Sentiment]} % Placeholder Box
    \caption{Average text length (characters) per sentiment class.}
    \label{fig:avg_length}
\end{figure}

\subsection{Evaluation Metrics}
\label{subsec:exp_metrics}

To evaluate and compare the performance of the three models, we use standard classification metrics derived from the confusion matrix on the test set:
\begin{itemize}[itemsep=0pt, topsep=3pt]
    \item \textbf{Accuracy}: The overall proportion of correctly classified tweets.
    $$ \text{Accuracy} = \frac{\text{Number of Correct Predictions}}{\text{Total Number of Predictions}} $$
    While simple, accuracy can be misleading on imbalanced datasets.
    \item \textbf{Precision}: For a given class, the proportion of tweets predicted as that class that actually belong to it. High precision indicates few false positives.
    $$ \text{Precision} = \frac{\text{True Positives (TP)}}{\text{TP} + \text{False Positives (FP)}} $$
    \item \textbf{Recall (Sensitivity)}: For a given class, the proportion of tweets actually belonging to that class that were correctly identified. High recall indicates few false negatives.
    $$ \text{Recall} = \frac{\text{TP}}{\text{TP} + \text{False Negatives (FN)}} $$
    \item \textbf{F1-Score}: The harmonic mean of Precision and Recall, providing a single balanced measure for each class.
    $$ \text{F1-Score} = 2 \times \frac{\text{Precision} \times \text{Recall}}{\text{Precision} + \text{Recall}} $$
\end{itemize}
We report these metrics per class, as well as macro-averaged and weighted-averaged scores. Macro-average calculates the metric independently for each class and then averages (treating all classes equally), while weighted-average accounts for class imbalance by weighting the average by the number of true instances for each class (support).

\subsection{Experiment 1: Lexicon-Based (TextBlob) Results}
\label{subsec:exp_lexicon}

The TextBlob model was applied directly to the test set texts (specifically, experiments were run using both `text` and `selected_text`, with results reported for analysis on `text` for direct comparison with other models). No training is involved. % Escaped underscore

\textbf{Results:} [Update with your actual results]
\begin{itemize}[itemsep=0pt, topsep=3pt]
    \item Overall Accuracy: Approx. 0.60
    \item Macro Average F1-Score: Approx. 0.57
    \item Performance varied significantly by class, often performing better on positive/neutral than negative.
\end{itemize}

\begin{figure}[H]
    \centering
    % \includegraphics[width=\columnwidth]{plots/model_analysis/lexicon_polarity_vs_subjectivity.png} % Replace placeholder
    \framebox[0.9\columnwidth][c]{[Figure Placeholder: TextBlob Polarity vs Subjectivity]} % Placeholder Box
    \caption{Polarity vs. Subjectivity scores from TextBlob for test set tweets, colored by true sentiment.}
    \label{fig:lex_polarity_subj}
\end{figure}

\begin{figure}[H]
    \centering
    % \includegraphics[width=\columnwidth]{plots/model_analysis/lexicon_polarity_distribution.png} % Replace placeholder
    \framebox[0.9\columnwidth][c]{[Figure Placeholder: TextBlob Polarity Distribution]} % Placeholder Box
    \caption{Distribution of TextBlob polarity scores for each true sentiment class.}
    \label{fig:lex_polarity_dist}
\end{figure}

\textbf{Discussion:} The TextBlob approach is extremely fast but yields the lowest accuracy among the three models. Figure \ref{fig:lex_polarity_subj} shows considerable overlap in polarity/subjectivity scores between the true classes, especially between neutral and negative sentiments. Figure \ref{fig:lex_polarity_dist} confirms that while positive tweets often receive positive polarity scores, negative and neutral tweets frequently fall into similar, near-zero polarity ranges, making discrimination based on simple thresholds difficult. The model's performance is highly dependent on its internal lexicon matching the vocabulary and sentiment expression style of the tweets.

\subsection{Experiment 2: SVM Results}
\label{subsec:exp_svm}

The LinearSVC model was trained on the TF-IDF vectors (unigrams and bigrams, max 5000 features) of the 80\% training split and evaluated on the 20\% test split.

\textbf{Results:} [Update with your actual results]
\begin{itemize}[itemsep=0pt, topsep=3pt]
    \item Overall Accuracy: Approx. 0.67
    \item Macro Average F1-Score: Approx. 0.67
    \item Balanced performance across classes, slightly better than Lexicon.
\end{itemize}

\begin{figure}[H]
    \centering
    % \includegraphics[width=0.8\columnwidth]{plots/svm_confusion_matrix.png} % Example path
    \framebox[0.8\columnwidth][c]{[Figure Placeholder: SVM Confusion Matrix]} % Placeholder Box
    \caption{Confusion Matrix for the SVM classifier on the test set.}
    \label{fig:svm_cm}
\end{figure}

\begin{figure}[H]
    \centering
    % \includegraphics[width=\columnwidth]{plots/svm_actual_vs_predicted_dist.png} % Example path
    \framebox[0.9\columnwidth][c]{[Figure Placeholder: SVM Actual vs Predicted Distribution]} % Placeholder Box
    \caption{Comparison of actual and SVM-predicted sentiment distributions on the test set.}
    \label{fig:svm_dist_comp}
\end{figure}

\begin{figure}[H]
    \centering
    % \includegraphics[width=\columnwidth]{plots/svm_per_class_metrics.png} % Example path
     \framebox[0.9\columnwidth][c]{[Figure Placeholder: SVM Per-Class Metrics (P, R, F1)]} % Placeholder Box
   \caption{Precision, Recall, and F1-Score per class for the SVM model.}
    \label{fig:svm_prf1}
\end{figure}

\begin{figure}[H]
    \centering
    % \includegraphics[width=\columnwidth]{plots/svm_text_length_correctness.png} % Example path
     \framebox[0.9\columnwidth][c]{[Figure Placeholder: SVM Text Length Correctness]} % Placeholder Box
    \caption{Distribution of text length for correctly vs incorrectly classified tweets by SVM.}
    \label{fig:svm_len_correct}
\end{figure}

\begin{figure}[H]
    \centering
    % \includegraphics[width=\columnwidth]{plots/model_analysis/svm_top_features_per_class.png} % Example path
     \framebox[0.9\columnwidth][c]{[Figure Placeholder: SVM Top N-grams per Class]} % Placeholder Box
    \caption{Most influential N-grams identified by the LinearSVC model for each sentiment class.}
    \label{fig:svm_features}
\end{figure}


\textbf{Discussion:} The SVM classifier significantly outperforms the lexicon-based approach, achieving better balance across classes (Figure \ref{fig:svm_prf1}). The confusion matrix (Figure \ref{fig:svm_cm}) reveals common confusions, often between neutral and positive/negative classes. The comparison of actual vs predicted distributions (Figure \ref{fig:svm_dist_comp}) shows [Describe observation, e.g., if SVM over/under predicts a class]. Analysis of text length versus prediction correctness (Figure \ref{fig:svm_len_correct}) suggests [Describe observation, e.g., SVM performs similarly regardless of length, or struggles with very short tweets]. The top features identified (Figure \ref{fig:svm_features}) provide some interpretability, highlighting words like 'happy', 'thanks' for positive and 'sad', 'sorry' for negative, which aligns with expectations. The performance evolution plot (not shown here, but generated by the script) indicated that performance generally improved with more training data but started to plateau, suggesting the model might be reaching its limit with the current features/data size.

\subsection{Experiment 3: LSTM Results}
\label{subsec:exp_lstm}

The LSTM model was trained specifically on positive and negative tweets from the training set for 10 epochs (or specify actual epochs used) and evaluated on the corresponding positive/negative tweets from the test set. *[Adapt this description if you trained a 3-class LSTM instead]*.

\textbf{Results:} [Update with your actual results - remember these are likely for 2 classes only if using the script as provided]
\begin{itemize}[itemsep=0pt, topsep=3pt]
    \item Accuracy (on Positive/Negative test set): Approx. 0.70 [Adjust based on your 2-class or 3-class results]
    \item Macro Average F1-Score (Positive/Negative): Approx. 0.70 [Adjust]
    \item Training took significantly longer than SVM.
\end{itemize}

\begin{figure}[H]
    \centering
    % \includegraphics[width=0.9\columnwidth]{results/lstm/lstm_training_history.png} % Example path
     \framebox[0.9\columnwidth][c]{[Figure Placeholder: LSTM Training History (Loss/Accuracy)]} % Placeholder Box
    \caption{LSTM training and validation loss/accuracy over epochs.}
    \label{fig:lstm_history}
\end{figure}

\begin{figure}[H]
    \centering
    % \includegraphics[width=\columnwidth]{plots/model_analysis/lstm_misclassified_wordcloud.png} % Example path
     \framebox[0.9\columnwidth][c]{[Figure Placeholder: LSTM Misclassified Word Cloud]} % Placeholder Box
    \caption{Frequent words in tweets misclassified by the LSTM model.}
    \label{fig:lstm_wc_misclassified}
\end{figure}

\begin{figure}[H]
    \centering
    % \includegraphics[width=\columnwidth]{plots/model_analysis/lstm_correctly_classified_wordcloud.png} % Example path
     \framebox[0.9\columnwidth][c]{[Figure Placeholder: LSTM Correctly Classified Word Cloud]} % Placeholder Box
    \caption{Frequent words in tweets correctly classified by the LSTM model.}
    \label{fig:lstm_wc_correct}
\end{figure}


\textbf{Discussion:} The LSTM model, even when trained only on binary classification (positive vs. negative) and for a limited number of epochs, showed slightly better performance metrics compared to the 3-class SVM on the subset of data it was trained on. [If you ran a 3-class LSTM, compare its 3-class results directly to SVM/Lexicon]. The training history (Figure \ref{fig:lstm_history}) shows [Describe trend, e.g., convergence, signs of overfitting/underfitting]. The word clouds (Figures \ref{fig:lstm_wc_misclassified} and \ref{fig:lstm_wc_correct}) offer qualitative insights. Words appearing frequently in misclassified tweets might include [Mention examples, e.g., negations, potentially sarcastic terms, or neutral words appearing in affectively charged contexts], suggesting areas where the model struggles with nuances. Conversely, correctly classified tweets often contain clearer sentiment indicators. The LSTM's ability to model sequences likely contributes to its performance, but it comes at the cost of higher computational requirements and complexity compared to SVM. Further hyperparameter tuning and more training epochs could potentially yield better results.

\subsection{Overall Comparison and Aggregation Analysis}
\label{subsec:exp_comparison} % Changed label to avoid duplicate

Table \ref{tab:comparison} summarizes the overall performance metrics (Macro Average F1-Score and Accuracy) for the three primary models evaluated on the 3-class task (Note: LSTM results need adaptation or separate reporting if only run on 2 classes).

\begin{table}[H]
\centering
\caption{Overall Performance Comparison (3-Class Evaluation)}
\label{tab:comparison}
\begin{tabular}{@{}lcc@{}}
\toprule
\textbf{Model} & \textbf{Accuracy} & \textbf{Macro Avg F1-Score} \\
\midrule
SVM        & [Your SVM Accuracy] & [Your SVM Macro F1] \\ % e.g., 0.67 & 0.67
Lexicon    & [Your Lexicon Accuracy] & [Your Lexicon Macro F1] \\ % e.g., 0.60 & 0.57
LSTM       & [Your 3-Class LSTM Acc]* & [Your 3-Class LSTM F1]* \\ % e.g., ~0.70* & ~0.70*
% Naive Bayes& [Your NB Accuracy] & [Your NB Macro F1] \\ % Add if you included NB comparison
\bottomrule
\end{tabular}
\caption*{\footnotesize{*LSTM results shown are indicative and may require re-running the model for 3-class evaluation.}}
\end{table}

\begin{figure}[H]
    \centering
    % \includegraphics[width=\columnwidth]{plots/model_comparison_prf1.png} % Example path
     \framebox[0.9\columnwidth][c]{[Figure Placeholder: Model Comparison (Precision, Recall, F1)]} % Placeholder Box
    \caption{Detailed comparison of Macro Average Precision, Recall, and F1-Score across models.}
    \label{fig:prf1_comparison}
\end{figure}

The results indicate that [Summarize findings, e.g., the deep learning approach (LSTM) potentially offers the highest accuracy, followed by SVM, with the Lexicon method performing worst]. Figure \ref{fig:prf1_comparison} provides a more detailed view, showing the trade-offs in precision and recall for each model. [Discuss the PRF1 plot results - e.g., SVM might have balanced P/R, Lexicon might have higher P but lower R].

Additionally, the MongoDB aggregation analysis provided insights into the dataset structure (Figure \ref{fig:sentiment_dist} and \ref{fig:avg_length}), confirming the class distribution and average text lengths efficiently without loading the entire dataset into memory for these specific calculations. This showcases an advantage of using a database with server-side processing capabilities.

\section{Conclusion}
\label{sec:conclusion}

This project presented a comparative analysis of lexicon-based, SVM, and LSTM methods for sentiment classification of tweets, using data managed in MongoDB Atlas. Our experiments demonstrated the typical trade-offs between these approaches.

The LSTM model achieved the highest performance metrics [Adjust based on your final results, especially if comparing 3-class models], showcasing the power of deep learning for capturing sequence information in text. However, its complexity, training time, and data requirements are significant considerations. The SVM provided a strong baseline performance with much faster training, making it a viable option when computational resources are limited or datasets are smaller. The TextBlob lexicon-based method, while simple and fast, yielded the lowest accuracy, highlighting the limitations of relying solely on predefined word scores for nuanced social media text.

Our use of MongoDB Atlas proved effective for managing the tweet data, offering flexibility and scalability. The ability to perform server-side aggregations was also demonstrated as an efficient way to derive dataset statistics.

**Limitations:** This study has several limitations. The LSTM model was trained for a limited number of epochs and potentially only on a binary subset; further tuning and training could improve results. Hyperparameter optimization was minimal for all models. The dataset size, while substantial, might still not be large enough to fully exploit the capabilities of deep learning models. We did not explicitly handle sarcasm or complex linguistic nuances beyond what the models learned implicitly. The evaluation focused on standard metrics and did not include qualitative error analysis in depth.

**Future Work:** Several avenues exist for future work:
\begin{itemize}[itemsep=0pt, topsep=3pt]
    \item **Sentiment Extraction:** Adapt the project to focus on extracting the specific text span (\texttt{selected\_text}) supporting the sentiment, which is a more challenging but relevant task for this dataset format. % Escaped underscore
    \item **Transformer Models:** Implement and compare state-of-the-art Transformer models like BERT or RoBERTa, potentially fine-tuning pre-trained versions for this task.
    \item **Hyperparameter Optimization:** Conduct systematic hyperparameter tuning for SVM and LSTM using techniques like Grid Search or Randomized Search.
    \item **Feature Engineering:** Explore more advanced features for SVM, such as incorporating sentiment lexicons or word embeddings.
    \item **Error Analysis:** Perform a more detailed qualitative analysis of misclassified tweets for each model to better understand their failure modes.
    \item **Scalability Testing:** Evaluate the performance and data handling capabilities on a significantly larger dataset to better assess the benefits of MongoDB and the scalability of the models.
    \item **Privacy Enhancements:** Implement the proposed PII removal techniques before analysis and evaluate its impact on model performance.
\end{itemize}

In conclusion, while deep learning models like LSTMs show promise for tweet sentiment analysis, classical methods like SVM remain competitive, especially considering implementation complexity and resource constraints. The choice of method should depend on the specific application requirements, available data, and computational budget.

% --- REFERENCES ---
% Use a proper bibliography style (e.g., bibtex) for a real report
% For now, using a simple list based on user's draft
\section*{References}
\label{sec:references}
\begin{thebibliography}{99} % Increased number to accommodate more refs if needed

\bibitem{ref:ahmad2017} % SVM paper placeholder
Ahmad, W., et al. "Sentiment Analysis of Tweets using SVM." \textit{International Journal of Computer Applications (IJCA)}, 2017. [Provide more complete citation details]

\bibitem{ref:roy2024} % LSTM/Hate speech paper placeholder
Roy, S. S., Roy, A., Samui, P., Gandomi, M., & Gandomi, A. H. "Hateful Sentiment Detection in Real-Time Tweets: An LSTM-Based Comparative Approach." \textit{IEEE Transactions on Computational Social Systems}, 11(4), 5028-5037, 2024.

\bibitem{ref:catelli2023} % Lexicon paper placeholder
Catelli, R., et al. "Lexicon-based sentiment analysis for COVID-19 related discussion in Italian tweets." \textit{Computers in Biology and Medicine}, 153, 106514, 2023.

% --- Add other references mentioned ---

\bibitem{ref:feldman2013}
Feldman, R. "Techniques and applications for sentiment analysis." \textit{Communications of the ACM}, 56(4), 82-89, 2013.

\bibitem{ref:sentiwordnet}
Baccianella, S., Esuli, A., & Sebastiani, F. "SentiWordNet 3.0: An Enhanced Lexical Resource for Sentiment Analysis and Opinion Mining." \textit{Proceedings of LREC}, 2010.

\bibitem{ref:vader}
Hutto, C. J., & Gilbert, E. "VADER: A Parsimonious Rule-based Model for Sentiment Analysis of Social Media Text." \textit{Proceedings of the Eighth International AAAI Conference on Weblogs and Social Media (ICWSM-14)}, 2014.

\bibitem{ref:textblob}
Loria, S., et al. TextBlob: Simplified Text Processing. \url{https://textblob.readthedocs.io/}. (Accessed: [Insert Date])

\bibitem{ref:pang_lee_sentiment}
Pang, B., & Lee, L. "Opinion mining and sentiment analysis." \textit{Foundations and Trends® in Information Retrieval}, 2(1–2), 1-135, 2008.

\bibitem{ref:svm_cortes}
Cortes, C., & Vapnik, V. "Support-vector networks." \textit{Machine Learning}, 20(3), 273-297, 1995.

\bibitem{ref:lstm_hochreiter}
Hochreiter, S., & Schmidhuber, J. "Long short-term memory." \textit{Neural computation}, 9(8), 1735-1780, 1997.

\bibitem{ref:gru_cho}
Cho, K., et al. "Learning phrase representations using RNN encoder-decoder for statistical machine translation." \textit{arXiv preprint arXiv:1406.1078}, 2014.

\bibitem{ref:cnn_kim}
Kim, Y. "Convolutional neural networks for sentence classification." \textit{arXiv preprint arXiv:1408.5882}, 2014.

\bibitem{ref:attention_bahdanau}
Bahdanau, D., Cho, K., & Bengio, Y. "Neural machine translation by jointly learning to align and translate." \textit{arXiv preprint arXiv:1409.0473}, 2014.

\bibitem{ref:transformer_vaswani}
Vaswani, A., et al. "Attention is all you need." \textit{Advances in neural information processing systems}, 30, 2017.

\bibitem{ref:bert_devlin}
Devlin, J., Chang, M. W., Lee, K., & Toutanova, K. "Bert: Pre-training of deep bidirectional transformers for language understanding." \textit{arXiv preprint arXiv:1810.04805}, 2018.

\bibitem{ref:word2vec}
Mikolov, T., et al. "Efficient Estimation of Word Representations in Vector Space." \textit{arXiv preprint arXiv:1301.3781}, 2013.

\bibitem{ref:glove}
Pennington, J., Socher, R., & Manning, C. D. "Glove: Global vectors for word representation." \textit{Proceedings of the 2014 conference on empirical methods in natural language processing (EMNLP)}, 2014.

\end{thebibliography}

\end{document}